\chapter*{How to use this book}
\addcontentsline{toc}{chapter}{How to use this book}
\markboth{How to use this book}{How to use this book}

\section*{The shape of the book}
The 52 chapters are grouped into twelve parts. The parts form a dependency structure
rather than a single line. The foundations in Part~II feed both the biological
abstractions of Part~III and the two model directions, DOGMA (Part~IV) and Hermon DNA
(Part~V). Part~VI's experimental contract governs every later comparison. The two engines
(Parts~VII and~VIII) are developed separately and meet only in Evolutor's runtime and
compiler (Part~IX). Industrial engineering, applications and synthesis then build on
that layer. Figure~\ref{book:readingmap} shows the structure and five suggested reading
routes.

\BackPlate{reading-map}{The twelve parts and their main dependencies. Solid arrows mean
``is needed by''. The coloured markers show five reading routes. Each route lists the
parts it reads in order and can skip the others on a first pass. Chapter-level
prerequisites are printed on every part page.}{book:readingmap}

\begin{center}\small
\begin{tabular}{@{}P{0.27\linewidth}P{0.33\linewidth}P{0.32\linewidth}@{}}
\toprule
Route & Chapters & What it emphasizes\\\midrule
Model builder & 2--10, then 18--29 & Objectives and state semantics before comparing architectures\\
Formal-systems reader & 11--17, then 39--43 & Separating biological analogy from typed semantics and legal lowering\\
Inference engineer & 23, 25--38 & Keeping scalar and full-model oracles while storage and scheduling change\\
Deployment engineer & 28--30, 32--38, 44--47 & Contracts that a reference test does not establish\\
Application researcher & 2, 5, 17, 21, 24--29, 48--52 & Datasets, interventions and negative evidence\\
\bottomrule
\end{tabular}
\end{center}
Chapter~1 is common to every route, and Chapter~52 together with the Epilogue closes each
one.

\section*{The anatomy of a chapter}
Every chapter follows the same progression, used with judgement rather than as a
template:
\begin{enumerate}
\item \textbf{A concrete problem and a prediction.} Each chapter opens with a small case
whose answer you can guess before reading on, such as ``must the resumed run take the same
next step?'' or ``which tokens change when one base changes?''.
\item \textbf{A derivation.} The mechanism is derived with every symbol, shape and unit
defined, and its premises stated next to the result.
\item \textbf{An implementation.} The printed listings are extracted from the tested
reference program at build time, so the code you read is the code that runs.
\item \textbf{A test or counterexample.} Each result is checked by an independent oracle,
by a deliberately broken negative control, or by a counterexample that shows a missing
premise matters.
\item \textbf{The consequence.} The chapter says what the result licenses, and what it
does not.
\end{enumerate}
The blue box at the start of each chapter lists its goals. The green \emph{Key ideas} box
before the exercises states what should stay with you. Exercises come with worked
solutions or marking rubrics. A closing \emph{Sources and evidence} note records which
primary sources were consulted and how deeply each was read.

\section*{Conventions}
\begin{itemize}
\setlength{\itemsep}{1pt}
\item Arrays and token positions are zero-based. Genomic intervals are half-open,
$[a,b)$, unless a file format states otherwise (Appendix~A).
\item Tensor shapes are written in square brackets, for example $[B,T,D]$ for batch,
time and width. Row vectors are used when the text follows PyTorch storage conventions.
\item $\log$ and $\ln$ denote natural logarithms, and losses are in nats unless bits are
stated. MiB and KiB are powers of 1024; MB and GB are powers of 1000.
\item DNA strands are written $5'$ to $3'$. The reverse complement $R(x)$ reverses the
order and complements each base; complementing alone is a different operation.
\item ``Tick'' means a logical step of a reference program, not elapsed time. Every
timing, throughput or memory figure states whether it was measured, modelled or
assigned.
\end{itemize}
Appendix~B collects the notation, and Appendix~C collects the evidence classes and the
claim-ledger template used throughout.

\section*{Running the companion code}
Each chapter from~4 onward has a reference program, generated tables and a test file. With
the pinned Python environment installed, the commands below rebuild the whole volume and
run every test. Rebuild the chapter assets before running the tests on a fresh checkout,
because several tests compare printed tables and listings with regenerated ones.
\begin{Verbatim}[fontsize=\small]
python3 scripts/check-chapter-structure.py
python3 scripts/build-convergence.py --through 52 --refresh-assets
python3 -m pytest -q -o addopts=''
\end{Verbatim}
A single chapter can be rebuilt and checked with \texttt{python3 drafts/chNN/build.py
--check}. All examples run on a CPU. None needs a GPU, a network service or private data.
